\subsection{ライフサイクル}
\subsubsection{ライフサイクルの定義、プロセス分類、用語}
ライフサイクルとは、システム、製品、サービス、プロジェクトなどが、構想から廃止まで進化していく全体である。ソフトウェア工学では、単に必要なプロセスを列挙するだけでは複雑さを十分に説明できないため、複数のプロセスと制約を含むライフサイクルとして捉える。

開発とは、利害関係者のニーズに従ってソフトウェアシステムを構築または変更する重要な段階である。製品ライフサイクルは、構想、提供、成長、成熟、廃止へ進む製品の変化として説明できる。これはソフトウェア固有ではなく、製品一般に当てはまる考え方である。

ソフトウェアシステムには、単体テストの対象になり得る原子的なソフトウェア構成要素、すなわちソフトウェア単位が含まれる。システム全体のライフサイクルは、発想、取得、供給、開発、運用、進化、保守、廃止までを含む。

\paragraph{表4：ISO/IEC/IEEE 12207 に基づくプロセス分類の学習用整理}
\par\smallskip\noindent\refstepcounter{table}\label{tab:ch10-04}表 \thetable: 表4：ISO/IEC/IEEE 12207 に基づくプロセス分類の学習用整理における分類・役割・主な例\par\smallskip
表 \thetable は、表4：ISO/IEC/IEEE 12207 に基づくプロセス分類の学習用整理における分類・役割・主な例を比較・確認しやすい形で整理したものである。\par\smallskip
\begin{longtable}{>{\raggedright\arraybackslash}p{0.16\linewidth}>{\raggedright\arraybackslash}p{0.30\linewidth}>{\raggedright\arraybackslash}p{0.46\linewidth}}
\toprule
分類 & 役割 & 主な例\\
\midrule
\endfirsthead
\toprule
分類 & 役割 & 主な例\\
\midrule
\endhead
技術プロセス & ソフトウェア製品を作り、進化させ、運用し、廃止するための工学的活動である。 & 業務分析、利害関係者ニーズ定義、要求定義、アーキテクチャ、設計、実装、統合、検証、妥当性確認、移行、運用、保守、廃止。\\
技術管理プロセス & プロジェクトを計画し、制御し、リスクや品質を管理する活動である。 & プロジェクト計画、評価と制御、意思決定、リスク管理、構成管理、情報管理、測定、品質保証。\\
組織的プロジェクト支援プロセス & 組織としてプロジェクトを支える基盤を整える活動である。 & ライフサイクルモデル管理、インフラ管理、ポートフォリオ管理、人材管理、品質管理、知識管理。\\
合意プロセス & 取得・供給など、組織間または関係者間の合意を支える活動である。 & 取得プロセス、供給プロセス。\\
\bottomrule
\end{longtable}
\par\smallskip
\begin{center}
\centering
\IfFileExists{assets/generated_figures/ch10/fig-ch10-03.png}{\includegraphics[width=0.96\linewidth]{assets/generated_figures/ch10/fig-ch10-03.png}}{\fbox{\parbox[c][48mm][c]{0.9\linewidth}{\centering 画像生成待ち\\四つのプロセスカテゴリの図}}}
\par\smallskip
\refstepcounter{figure}\label{fig:ch10-03}図 \thefigure: 四つのプロセスカテゴリの図
\end{center}
図 \thefigure は、四つのプロセスカテゴリの図を視覚的に整理したものである。
\par\smallskip
\subsubsection{ライフサイクルが必要な理由}
ソフトウェア製品の作成、運用、廃止には、多数のプロセス、活動、タスク、制約が関わる。さらに、ソフトウェアシステムは、人、手順、ハードウェアを含むため、単純な工程表だけでは表現しにくい。

プロセスの入力は、多くの場合、別のプロセスの出力である。したがって、各プロセスは独立しているのではなく、相互に依存している。この相互依存が、ソフトウェア工学プロセス全体を複雑にする。

ライフサイクル仕様は、工学的な取り組みを可能にする強力な道具である。ライフサイクルを定義することは、各プロセスと制約を定義することであり、関係者間のコミュニケーション、技術管理、測定、評価、改善、品質管理の土台になる。

\begin{pointbox}{実務での見方}「どの工程を先にやるか」だけでなく、「どの活動の出力を誰が使うか」「何を測れば状態を把握できるか」「変更が起きたときにどこへ戻るか」を考えると、ライフサイクルを実務的に理解しやすい。\end{pointbox}
\subsubsection{プロセスモデルとライフサイクルモデルの概念}
プロセスモデルは、活動、入力、出力、制約、関係をどのように表すかというモデルである。これに対して、ライフサイクルモデルは、活動をどの順序や反復構造で配置するかを示す型である。

プロジェクト固有の活動列は、標準で定義された活動を、選択したソフトウェアライフサイクルモデルに写像することで作られる。つまり、現実のプロジェクトのライフサイクルは、標準や組織規程をそのまま丸写しするものではなく、選んだモデルに従って具体化するものである。

代表的なライフサイクルモデルには、ウォーターフォールモデル、Vモデル、増分モデル、スパイラルモデル、アジャイルモデルがある。

\paragraph{表5：プロセスモデルとライフサイクルモデルの違い}
\par\smallskip\noindent\refstepcounter{table}\label{tab:ch10-05}表 \thetable: 表5：プロセスモデルとライフサイクルモデルの違いにおける観点・プロセスモデル・ライフサイクルモデル\par\smallskip
表 \thetable は、表5：プロセスモデルとライフサイクルモデルの違いにおける観点・プロセスモデル・ライフサイクルモデルを比較・確認しやすい形で整理したものである。\par\smallskip
\begin{longtable}{>{\raggedright\arraybackslash}p{0.16\linewidth}>{\raggedright\arraybackslash}p{0.30\linewidth}>{\raggedright\arraybackslash}p{0.46\linewidth}}
\toprule
観点 & プロセスモデル & ライフサイクルモデル\\
\midrule
\endfirsthead
\toprule
観点 & プロセスモデル & ライフサイクルモデル\\
\midrule
\endhead
主な関心 & 活動が何を入力とし、何を出力するか。 & 活動をどの順序、段階、反復で実施するか。\\
例 & IDEF0、BPMN、UML活動図、データフロー図。 & ウォーターフォール、Vモデル、増分、スパイラル、アジャイル。\\
実務での使い方 & 作業の中身と接続を説明する。 & プロジェクト全体の進み方を説明する。\\
\bottomrule
\end{longtable}
\subsubsection{開発ライフサイクルモデルの考え方}
各ソフトウェアシステムには、利害関係者の業務的・技術的ニーズを反映した固有の特徴がある。したがって、適切なライフサイクルは、その特徴を考慮して決める必要がある。

\term{予測型ライフサイクル}{Predictive Life Cycle}では、範囲、時間、費用を早い段階で決める。これは、実装される要求集合が実質的に閉じており、大きな変更が起きないことを前提にする。

反復型ライフサイクルでは、範囲は比較的早く定まるが、製品理解が深まるにつれて時間や費用の見積りを更新する。進化型ライフサイクルでは、要求や顧客ニーズの変化、または要求の段階的導入により、製品やサービスが寿命を通じて変わる。増分型ライフサイクルでは、決められた期間の反復を通じて成果物へ機能を順次追加する。

継続的開発は、自動化ツールを用いて、新しいシステムやソフトウェアをステージング環境や各種テスト環境へ頻繁にリリースできるようにする実践群である。要求仕様の変更を許すライフサイクルは、変更に開かれているといえる。

\paragraph{表6：開発ライフサイクルの考え方の比較}
\par\smallskip\noindent\refstepcounter{table}\label{tab:ch10-06}表 \thetable: 表6：開発ライフサイクルの考え方の比較における考え方・特徴・初学者向けの見分け方\par\smallskip
表 \thetable は、表6：開発ライフサイクルの考え方の比較における考え方・特徴・初学者向けの見分け方を比較・確認しやすい形で整理したものである。\par\smallskip
\begin{longtable}{>{\raggedright\arraybackslash}p{0.16\linewidth}>{\raggedright\arraybackslash}p{0.30\linewidth}>{\raggedright\arraybackslash}p{0.46\linewidth}}
\toprule
考え方 & 特徴 & 初学者向けの見分け方\\
\midrule
\endfirsthead
\toprule
考え方 & 特徴 & 初学者向けの見分け方\\
\midrule
\endhead
予測型 & 範囲・時間・費用を早期に決め、要求の大幅変更を想定しにくい。 & 最初に決めた計画への適合を重視する。\\
反復型 & 同じ種類の活動を繰り返し、理解の進展に応じて見積りを更新する。 & 作りながら学び、次の反復へ反映する。\\
進化型 & 製品やサービスが寿命を通じて変化する。 & ニーズや技術の変化に合わせて育てる。\\
増分型 & 反復ごとに機能を追加し、最後に十分な能力を持つ成果物になる。 & 小さな成果を積み上げる。\\
継続的開発 & 自動化により、リリースやテスト環境への反映を高頻度に行う。 & 短い間隔で出し、測り、直す。\\
\bottomrule
\end{longtable}
\par\smallskip
\begin{center}
\centering
\IfFileExists{assets/generated_figures/ch10/fig-ch10-04.png}{\includegraphics[width=0.96\linewidth]{assets/generated_figures/ch10/fig-ch10-04.png}}{\fbox{\parbox[c][48mm][c]{0.9\linewidth}{\centering 画像生成待ち\\ライフサイクルモデル比較図}}}
\par\smallskip
\refstepcounter{figure}\label{fig:ch10-04}図 \thefigure: ライフサイクルモデル比較図
\end{center}
図 \thefigure は、ライフサイクルモデル比較図を視覚的に整理したものである。
\par\smallskip
\subsubsection{開発ライフサイクルモデルと工学的側面}
ウォーターフォールモデルは、ソフトウェア工学の初期に広く知られるようになった予測型のモデルである。要求、予備設計、詳細設計、コーディング、テストといったフェーズを順序立てて進め、前のフェーズが完了するまで次に進まない厳格な進め方を重視する。文書駆動型であり、ソフトウェア危機に対する初期の工学的対応として意味を持った。

Vモデルは、ウォーターフォール型の変種または拡張として理解できる。増分型ライフサイクルでは、工程が完全な直列ではなく重なり合いながら進む。増分型であっても、要求を先に閉じるなら予測型にもなり得る。

スパイラルモデルは、進化型でリスク駆動のモデルである。要求分析、設計、コーディング、統合、テストなどを反復し、リスクを見ながら完成へ近づける。迅速プロトタイピングは、早期に試作品を作り、早いフィードバックを得ることを重視する。統一プロセス、RUP、OpenUP は、反復型・増分型のプロセス枠組みとして知られる。

アジャイル宣言は、ソフトウェア工学コミュニティの考え方を大きく変えた。アジャイルは、要求が利用者ニーズの変化に応じてどの段階でも変更され得ること、チームと顧客の信頼、コミュニケーション、価値の継続的提供、技術的卓越性を重視する。重要なのは、アジャイルは文書を不要とするものではないという点である。必要な文書は必要である。

アジャイルに関する誤解も多い。アジャイルは単一の方法ではない。文書を作らないから速い、という意味でもない。エクストリームプログラミング、スクラム、スプリント計画、ふりかえり、ペアプログラミングなど、多様な方法と実践が存在する。大規模プロジェクトやポートフォリオへアジャイルを拡張することは、いまなお難しい課題である。

DevOps は、開発、運用、関連する利害関係者のコミュニケーションと協力を高め、ソフトウェアとシステムの構築、パッケージ化、デプロイ、運用、改善をつなげる考え方である。より頻繁なリリースを可能にすることは、適切なプロセス管理のもとでは競争力になる。

ウォーターフォール対アジャイルのような論争は、歴史的背景を持つ。ソフトウェア工学では、好みや流行ではなく、経験的測定に基づいて議論すべきである。プロセスや製品の測定は、意思決定に欠かせない。

\par\smallskip
\begin{center}
\centering
\IfFileExists{assets/generated_figures/ch10/fig-ch10-05.png}{\includegraphics[width=0.96\linewidth]{assets/generated_figures/ch10/fig-ch10-05.png}}{\fbox{\parbox[c][48mm][c]{0.9\linewidth}{\centering 画像生成待ち\\代表的ライフサイクルモデルの位置づけ図}}}
\par\smallskip
\refstepcounter{figure}\label{fig:ch10-05}図 \thefigure: 代表的ライフサイクルモデルの位置づけ図
\end{center}
図 \thefigure は、代表的ライフサイクルモデルの位置づけ図を視覚的に整理したものである。
\par\smallskip
\subsubsection{SLCPの管理}
SLCP はソフトウェアライフサイクルプロセスである。一般的なライフサイクルには、概念、開発、生産、利用、支援、廃止という段階がある。

概念段階では、利害関係者のニーズを識別し、概念を探索し、解決策を提案する。開発段階では、利用者ニーズを表す要求を洗練し、解決策を作り、システムを構築し、必要な検証と妥当性確認を行う。生産段階は対象システムの性質により範囲が異なるが、一般にはシステムの生産とテストを含む。利用段階では、システムが利用者ニーズを満たすために運用される。支援段階では、満足できる運用を実現するための支援行為を行う。廃止段階では、定められた手順に従ってシステムを処分する。

これらの段階は、必ず直列に進むわけではない。実際のライフサイクル仕様では、段階間の遷移を定義する。個別のプロジェクトでは、利害関係者にとって意味のある固有の段階が設定され、それらが一般的段階に対応づけられる。

\paragraph{表7：一般的なライフサイクル段階}
\par\smallskip\noindent\refstepcounter{table}\label{tab:ch10-07}表 \thetable: 表7：一般的なライフサイクル段階における段階・主な内容\par\smallskip
表 \thetable は、表7：一般的なライフサイクル段階における段階・主な内容を比較・確認しやすい形で整理したものである。\par\smallskip
\begin{longtable}{>{\raggedright\arraybackslash}p{0.28\linewidth}>{\raggedright\arraybackslash}p{0.64\linewidth}}
\toprule
段階 & 主な内容\\
\midrule
\endfirsthead
\toprule
段階 & 主な内容\\
\midrule
\endhead
概念 & ニーズを識別し、概念と解決策を探索する。\\
開発 & 要求を洗練し、解決策を作り、検証と妥当性確認を行う。\\
生産 & 対象に応じて、システムの生産とテストを行う。\\
利用 & システムを運用し、利用者ニーズを満たす。\\
支援 & 満足できる運用のための保守・支援行為を行う。\\
廃止 & 定められた手順でシステムを処分する。\\
\bottomrule
\end{longtable}
\par\smallskip
\begin{center}
\centering
\IfFileExists{assets/generated_figures/ch10/fig-ch10-06.png}{\includegraphics[width=0.96\linewidth]{assets/generated_figures/ch10/fig-ch10-06.png}}{\fbox{\parbox[c][48mm][c]{0.9\linewidth}{\centering 画像生成待ち\\ライフサイクル段階と遷移の図}}}
\par\smallskip
\refstepcounter{figure}\label{fig:ch10-06}図 \thefigure: ライフサイクル段階と遷移の図
\end{center}
図 \thefigure は、ライフサイクル段階と遷移の図を視覚的に整理したものである。
\par\smallskip
\subsubsection{ソフトウェア工学プロセス管理}
プロセス管理とは、製品を開発したりサービスを実施したりする作業を、方向づけ、制御し、調整することである。ソフトウェア工学プロセスは複数の管理階層で統制される。

最下層には、要求定義、設計、実装、検証、運用、保守といった技術プロセスがある。第二層には、プロジェクト計画やリスク管理などを含む技術管理レベルがある。第三層には、知識管理、ライフサイクルモデル管理、ポートフォリオ管理などを扱う経営または組織的支援レベルがある。

この階層を区別すると、問題の所在を見誤りにくくなる。実装が遅れている原因が、技術者の作業ではなく、プロジェクト計画、ツール統合、人材配置、組織知識の不足にあることもある。

\subsubsection{ソフトウェアライフサイクル適応}
各ソフトウェアシステムには固有の特徴がある。製品規模、複雑さ、利害関係者、規制、品質水準、運用環境、組織文化が異なれば、適切なライフサイクルも異なる。

ライフサイクル適応では、関連する特徴を識別し、適切な標準や組織内文書を選び、開発戦略やライフサイクルモデル、段階、プロセスを選択する。そして、その判断と根拠を文書化する。

ここで重要なのは、標準に書かれた名前をすべて保つ必要はなく、すべてのプロセスをそのまま含める必要もないという点である。標準は土台であり、現実のプロジェクトでは適用範囲と調整理由を明確にすることが重要である。小規模組織向けの ISO/IEC 29110 系列は、12207 から派生した例として理解できる。

\begin{pointbox}{適応の勘所}「標準に従う」とは、標準を機械的に写すことではない。何を採用し、何を省き、なぜそうしたかを説明できる状態にすることである。\end{pointbox}
\subsubsection{実務上の考慮}
ライフサイクルプロセスを定義するとは、技術プロセス、組織的プロセス、技術管理プロセス、合意プロセスを具体化することである。つまり、何を作るかだけでなく、誰が支え、どう管理し、どのように合意するかを定義する。

ソフトウェア工学は誕生以来、学び続けてきた分野である。製品の複雑さは増し続けているため、製品特性と外部特性の両方を考慮する必要がある。利害関係者の特徴も、ライフサイクル選択に影響する。

見積りと測定は不可欠である。ライフサイクルの中で誤った見積りや不確実な見積りを使うと、プロジェクト失敗につながる。ただし、正確な見積りは容易ではない。したがって、測定値には精度と不確実性を明示し、プロセスと製品の状態を正確に把握する必要がある。

現在の傾向として、継続的デリバリーが重視される。大きなプロセスを長期間進めながら途中成果物を出さないと、不確実性が増える。アジャイルの考え方は、途中で価値を示し、コミュニケーションを高め、不確実性を下げる点に貢献している。

\begin{pointbox}{論争より測定}予測型かアジャイルかという議論は、現実のプロセス測定と製品測定に基づけるべきである。測定が不確実なら、プロセス選択そのものを振り返る必要がある。\end{pointbox}
\subsubsection{ソフトウェアプロセス基盤、ツール、手法}
ソフトウェアプロセスを定義する記法には、自然言語、活動とタスクの一覧、データフロー図、状態図、IDEF0、ペトリネット、UML活動図、BPMN などがある。記法を選ぶ目的は、プロセスを人が理解し、合意し、管理し、改善できるようにすることである。

プロセス基盤には、プロセス定義を支援するツールだけでなく、テスト、構成管理、チケット管理など、個別プロセスを支援するツールが含まれる。重要なのは、ツールが単独で存在するのではなく、互いに統合されることである。

たとえば、あるコード単位がテストに合格したなら、その情報はテストツールだけに閉じていては不十分である。構成管理ツールやチーム全体の状態把握に連携されることで、次の活動が正しく判断できる。

統合されたツール群は、ソフトウェア工学環境と呼ばれることがある。1980年代から1990年代には CASE という語がよく使われたが、万能薬として過度に期待された面もあった。現在でも、バージョン管理、テスト、チケット管理などの自動化は、成功するライフサイクル実装に不可欠である。

\paragraph{表8：プロセス定義に使われる表現例}
\par\smallskip\noindent\refstepcounter{table}\label{tab:ch10-08}表 \thetable: 表8：プロセス定義に使われる表現例における表現・向いている用途\par\smallskip
表 \thetable は、表8：プロセス定義に使われる表現例における表現・向いている用途を比較・確認しやすい形で整理したものである。\par\smallskip
\begin{longtable}{>{\raggedright\arraybackslash}p{0.28\linewidth}>{\raggedright\arraybackslash}p{0.64\linewidth}}
\toprule
表現 & 向いている用途\\
\midrule
\endfirsthead
\toprule
表現 & 向いている用途\\
\midrule
\endhead
自然言語 & 関係者が読みやすい説明。曖昧さには注意が必要である。\\
活動・タスク一覧 & 作業を漏れなく列挙したい場合。\\
データフロー図 & 情報や成果物の流れを示したい場合。\\
状態図 & 状態と遷移を明確にしたい場合。\\
IDEF0 & 入力、制約、出力、支援資源を整理したい場合。\\
UML活動図 & 活動の分岐、合流、並行を示したい場合。\\
BPMN & 業務プロセスや組織横断の流れを示したい場合。\\
\bottomrule
\end{longtable}
\par\smallskip
\begin{center}
\centering
\IfFileExists{assets/generated_figures/ch10/fig-ch10-07.png}{\includegraphics[width=0.96\linewidth]{assets/generated_figures/ch10/fig-ch10-07.png}}{\fbox{\parbox[c][48mm][c]{0.9\linewidth}{\centering 画像生成待ち\\ツール連携としてのソフトウェア工学環境の図}}}
\par\smallskip
\refstepcounter{figure}\label{fig:ch10-07}図 \thefigure: ツール連携としてのソフトウェア工学環境の図
\end{center}
図 \thefigure は、ツール連携としてのソフトウェア工学環境の図を視覚的に整理したものである。
\par\smallskip
\subsubsection{プロセス監視とソフトウェア製品の関係}
開発者や管理者は、ソフトウェア工学プロセスの実行を監視し、プロセス目標が満たされているか、リスクが高まっていないかを評価する必要がある。これはプロセス評価の一部であり、ソフトウェア工学マネジメントとも関係する。

\term{経験的方法}{Empirical Method}は、プロセスの評価と改善だけでなく、製品の評価と改善も支える。プロセス実行の目的は、利害関係者のニーズを満たす製品を得ることである。したがって、プロセスだけを測っても不十分であり、製品の状態とあわせて見る必要がある。

たとえば、計画通りにレビュー回数が実施されていても、欠陥が減っていないなら、プロセスは目的を果たしていない可能性がある。逆に、製品品質が高く見えても、チームが過負荷で属人的な作業に依存しているなら、長期的なプロセスリスクがある。

\par\smallskip
\begin{center}
\centering
\IfFileExists{assets/generated_figures/ch10/fig-ch10-08.png}{\includegraphics[width=0.96\linewidth]{assets/generated_figures/ch10/fig-ch10-08.png}}{\fbox{\parbox[c][48mm][c]{0.9\linewidth}{\centering 画像生成待ち\\プロセス監視と製品評価の関係図}}}
\par\smallskip
\refstepcounter{figure}\label{fig:ch10-08}図 \thefigure: プロセス監視と製品評価の関係図
\end{center}
図 \thefigure は、プロセス監視と製品評価の関係図を視覚的に整理したものである。
\par\smallskip
